\documentclass[10pt,a4paper]{amsart}
\usepackage{amsmath,amssymb,amsthm,graphicx}
\usepackage{verbatim}
\newenvironment{reference}{\comment}{\endcomment}
%\newenvironment{reference}{}{}
\newenvironment{slogan}{\comment}{\endcomment}
\newenvironment{history}{\comment}{\endcomment}

% For commutative diagrams we use Xy-pic
\usepackage[all]{xy}

% We use 2cell for 2-commutative diagrams.
\xyoption{2cell}
% Theorem environments.
%
\theoremstyle{plain}
\newtheorem{theorem}[subsection]{Theorem}
\newtheorem{proposition}[subsection]{Proposition}
\newtheorem{lemma}[subsection]{Lemma}

\theoremstyle{definition}
\newtheorem{definition}[subsection]{Definition}
\newtheorem{example}[subsection]{Example}
\newtheorem{exercise}[subsection]{Exercise}
\newtheorem{situation}[subsection]{Situation}

\theoremstyle{remark}
\newtheorem{remark}[subsection]{Remark}
\newtheorem{remarks}[subsection]{Remarks}

\numberwithin{equation}{subsection}

% Macros
%
\def\lim{\mathop{\mathrm{lim}}\nolimits}
\def\colim{\mathop{\mathrm{colim}}\nolimits}
\def\Spec{\mathop{\mathrm{Spec}}}
\def\Hom{\mathop{\mathrm{Hom}}\nolimits}
\def\Ext{\mathop{\mathrm{Ext}}\nolimits}
\def\SheafHom{\mathop{\mathcal{H}\!\mathit{om}}\nolimits}
\def\SheafExt{\mathop{\mathcal{E}\!\mathit{xt}}\nolimits}
\def\Sch{\mathit{Sch}}
\def\Mor{\mathop{\mathrm{Mor}}\nolimits}
\def\Ob{\mathop{\mathrm{Ob}}\nolimits}
\def\Sh{\mathop{\mathit{Sh}}\nolimits}
\def\NL{\mathop{N\!L}\nolimits}
\def\CH{\mathop{\mathrm{CH}}\nolimits}
\def\proetale{{pro\text{-}\acute{e}tale}}
\def\etale{{\acute{e}tale}}
\def\QCoh{\mathit{QCoh}}
\def\Ker{\mathop{\mathrm{Ker}}}
\def\Im{\mathop{\mathrm{Im}}}
\def\Coker{\mathop{\mathrm{Coker}}}
\def\Coim{\mathop{\mathrm{Coim}}}

% Boxtimes
%
\DeclareMathSymbol{\boxtimes}{\mathbin}{AMSa}{"02}

%
% Macros for moduli stacks/spaces
%
\def\QCohstack{\mathcal{QC}\!\mathit{oh}}
\def\Cohstack{\mathcal{C}\!\mathit{oh}}
\def\Spacesstack{\mathcal{S}\!\mathit{paces}}
\def\Quotfunctor{\mathrm{Quot}}
\def\Hilbfunctor{\mathrm{Hilb}}
\def\Curvesstack{\mathcal{C}\!\mathit{urves}}
\def\Polarizedstack{\mathcal{P}\!\mathit{olarized}}
\def\Complexesstack{\mathcal{C}\!\mathit{omplexes}}
% \Pic is the operator that assigns to X its picard group, usage \Pic(X)
% \Picardstack_{X/B} denotes the Picard stack of X over B
% \Picardfunctor_{X/B} denotes the Picard functor of X over B
\def\Pic{\mathop{\mathrm{Pic}}\nolimits}
\def\Picardstack{\mathcal{P}\!\mathit{ic}}
\def\Picardfunctor{\mathrm{Pic}}
\def\Deformationcategory{\mathcal{D}\!\mathit{ef}}
\usepackage{cleveref}
\crefname{originallemma}{original lemma}{original lemmas}
\Crefname{originallemma}{Original lemma}{Original lemmas}
\makeatletter
\newlabel{coherent-lemma-proper-finite-fibre-finite-in-neighbourhood}{{02OH}{source}{}{sourcepointer-1}{}}
\newlabel{coherent-lemma-proper-finite-fibre-finite-in-neighbourhood@cref}{{[originallemma][0][]02OH}{[1][0][]source}{}{}{}}
\makeatother
\crefname{originallemma}{original lemma}{original lemmas}
\Crefname{originallemma}{Original lemma}{Original lemmas}
\makeatletter
\newlabel{more-morphisms-lemma-proper-finite-fibre-finite-in-neighbourhood}{{02UP}{source}{}{sourcepointer-2}{}}
\newlabel{more-morphisms-lemma-proper-finite-fibre-finite-in-neighbourhood@cref}{{[originallemma][0][]02UP}{[1][0][]source}{}{}{}}
\makeatother
\crefname{originallemma}{original lemma}{original lemmas}
\Crefname{originallemma}{Original lemma}{Original lemmas}
\makeatletter
\newlabel{morphisms-lemma-image-proper-scheme-closed}{{01W6}{source}{}{sourcepointer-3}{}}
\newlabel{morphisms-lemma-image-proper-scheme-closed@cref}{{[originallemma][0][]01W6}{[1][0][]source}{}{}{}}
\makeatother
\title{A morphism from a proper dimension-1 scheme over a field is finite when no component maps to a single point}
\author{The Stacks Project Authors}
\date{stacks-project:curves.tex:local-snapshot-20261012; GNU FDL 1.2 or later}
\begin{document}
\maketitle
\noindent\textit{Editorial note.} All delivered mathematical text is the author's own, copied from the pinned Stacks Project chapter file sources/stacks-curves/pinned-original-source.tex; no statement or proof step was rewritten, reordered or removed.\par\medskip
\noindent\textit{Editorial note.} The counted proof is the author's complete argument: the properness of f, the closedness of f(X) and of point images, the fibre analysis showing fibres over closed points are finite because they avoid the one-dimensional generic points, the local finiteness from the proper-finite-fibre lemma in both its general and Noetherian forms, and the covering conclusion are all retained verbatim.\par\medskip
\noindent\textit{Editorial note.} The label lemma-finite names the counted statement, and the proof environment immediately following it in the source is delivered verbatim as its proof.\par\medskip
\noindent\textit{Editorial note.} The counted statement is carried by the author's own lemma environment in the chapter source.\par\medskip
\noindent\textit{Allowed standard prerequisites (editorial).}\begin{itemize}
\item Properness of morphisms from proper schemes to separated schemes and closedness of images. Morphisms, Lemma morphisms-lemma-image-proper-scheme-closed. Gives that f itself is proper, hence f(X) is closed and images of closed points are closed, the starting point of the fibre analysis.
\item Proper morphisms with finite fibres are finite locally on the base. More on Morphisms, Lemma more-morphisms-lemma-proper-finite-fibre-finite-in-neighbourhood; Cohomology of Schemes, Lemma coherent-lemma-proper-finite-fibre-finite-in-neighbourhood. Turns the pointwise finiteness of the closed fibres into finiteness of f over an open neighbourhood of each image point.
\end{itemize}\medskip
\ifdefined\hypertarget\hypertarget{sourcepointer-1}{}\fi
\noindent\textit{Source reference.} Original lemma 02OH: sources/stacks-chow/tags.txt, source lines 3516--3516. Prerequisite pointer invoked at proof line(s) 349; it names a standard fact of a sibling chapter that the counted argument uses, and every step of the counted proof built on it is retained verbatim in the delivered body.\par\medskip
\ifdefined\hypertarget\hypertarget{sourcepointer-2}{}\fi
\noindent\textit{Source reference.} Original lemma 02UP: sources/stacks-chow/tags.txt, source lines 3740--3740. Prerequisite pointer invoked at proof line(s) 346; it names a standard fact of a sibling chapter that the counted argument uses, and every step of the counted proof built on it is retained verbatim in the delivered body.\par\medskip
\ifdefined\hypertarget\hypertarget{sourcepointer-3}{}\fi
\noindent\textit{Source reference.} Original lemma 01W6: sources/stacks-chow/tags.txt, source lines 2471--2471. Prerequisite pointer invoked at proof line(s) 337; it names a standard fact of a sibling chapter that the counted argument uses, and every step of the counted proof built on it is retained verbatim in the delivered body.\par\medskip
\begin{lemma}
\label{lemma-finite}
Let $k$ be a field. Let $f : X \to Y$ be a morphism of
schemes over $k$. Assume
\begin{enumerate}
\item $Y$ is separated over $k$,
\item $X$ is proper of dimension $\leq 1$ over $k$,
\item $f(Z)$ has at least two points for every irreducible
component $Z \subset X$ of dimension $1$.
\end{enumerate}
Then $f$ is finite.
\end{lemma}

\begin{proof}
The morphism $f$ is proper by
Morphisms, Lemma \ref{morphisms-lemma-image-proper-scheme-closed}.
Thus $f(X)$ is closed and images of closed points are closed.
Let $y \in Y$ be the image of a closed point in $X$.
Then $f^{-1}(\{y\})$ is a closed subset of $X$ not
containing any of the generic points of irreducible components
of dimension $1$ by condition (3). It follows that $f^{-1}(\{y\})$
is finite. Hence $f$ is finite over an open neighbourhood of $y$
by
More on Morphisms, Lemma
\ref{more-morphisms-lemma-proper-finite-fibre-finite-in-neighbourhood}
(if $Y$ is Noetherian, then you can use the easier
Cohomology of Schemes, Lemma
\ref{coherent-lemma-proper-finite-fibre-finite-in-neighbourhood}).
Since we've seen above that there are enough of these points
$y$, the proof is complete.
\end{proof}

\end{document}
